\documentclass[11pt]{article}
\newcommand{\videoid}{EZlgC_DT7D0}
% Preámbulo compartido de las notas del curso de topología algebraica.
% Cada nota empieza con \documentclass[11pt]{article}, define \videoid{...}
% (el identificador del video de YouTube) y luego hace % Preámbulo compartido de las notas del curso de topología algebraica.
% Cada nota empieza con \documentclass[11pt]{article}, define \videoid{...}
% (el identificador del video de YouTube) y luego hace % Preámbulo compartido de las notas del curso de topología algebraica.
% Cada nota empieza con \documentclass[11pt]{article}, define \videoid{...}
% (el identificador del video de YouTube) y luego hace \input{preambulo}.
% Compilar con LuaLaTeX:  lualatex nombre-de-la-nota.tex
\usepackage[spanish,provide=*]{babel}
\usepackage[T1]{fontenc}
\usepackage{lmodern}
\usepackage[a4paper,margin=2.6cm]{geometry}
\usepackage{amsmath,amssymb,amsthm}
\usepackage{xcolor}
\usepackage{tikz}
\usetikzlibrary{arrows.meta,babel}
\usepackage[colorlinks=true,urlcolor=blue!50!black,linkcolor=blue!50!black]{hyperref}

\hyphenpenalty=10000 \exhyphenpenalty=10000 \emergencystretch=3em

% Colores de las figuras (los mismos que usa la página)
\definecolor{extremo}{RGB}{47,107,79}
\definecolor{azul}{RGB}{44,74,154}

% Definiciones numeradas por unidad
\theoremstyle{definition}
\newtheorem{defn}{Definición}
% Numeración por unidad: \unidad{u}{n} (después de \input{preambulo}) hace que las
% definiciones de la nota se llamen u.n, u.(n+1), ...; n es el primer número
% libre de esa unidad (ver la tabla de CONVENCIONES.md).
\newcommand{\unidad}[2]{\renewcommand{\thedefn}{#1.\arabic{defn}}\setcounter{defn}{\numexpr#2-1\relax}}

% Observaciones y erratas del video. Usar SOLO si hay algo que señalar:
%   \begin{erratas}
%     \errata{203}{3:23}{Se dice ..., pero debe decir ...}
%   \end{erratas}
% El primer argumento son los segundos (enlaza al minuto exacto del video),
% el segundo es la marca visible, el tercero es el texto.
\newenvironment{erratas}{\section*{Observaciones y erratas del video}\begin{itemize}}{\end{itemize}}
\newcommand{\errata}[3]{%
  \edef\errurl{https://www.youtube.com/watch?v=\videoid\string&t=#1}%
  \item[\textbf{\expandafter\href\expandafter{\errurl}{#2}}] #3}
.
% Compilar con LuaLaTeX:  lualatex nombre-de-la-nota.tex
\usepackage[spanish,provide=*]{babel}
\usepackage[T1]{fontenc}
\usepackage{lmodern}
\usepackage[a4paper,margin=2.6cm]{geometry}
\usepackage{amsmath,amssymb,amsthm}
\usepackage{xcolor}
\usepackage{tikz}
\usetikzlibrary{arrows.meta,babel}
\usepackage[colorlinks=true,urlcolor=blue!50!black,linkcolor=blue!50!black]{hyperref}

\hyphenpenalty=10000 \exhyphenpenalty=10000 \emergencystretch=3em

% Colores de las figuras (los mismos que usa la página)
\definecolor{extremo}{RGB}{47,107,79}
\definecolor{azul}{RGB}{44,74,154}

% Definiciones numeradas por unidad
\theoremstyle{definition}
\newtheorem{defn}{Definición}
% Numeración por unidad: \unidad{u}{n} (después de % Preámbulo compartido de las notas del curso de topología algebraica.
% Cada nota empieza con \documentclass[11pt]{article}, define \videoid{...}
% (el identificador del video de YouTube) y luego hace \input{preambulo}.
% Compilar con LuaLaTeX:  lualatex nombre-de-la-nota.tex
\usepackage[spanish,provide=*]{babel}
\usepackage[T1]{fontenc}
\usepackage{lmodern}
\usepackage[a4paper,margin=2.6cm]{geometry}
\usepackage{amsmath,amssymb,amsthm}
\usepackage{xcolor}
\usepackage{tikz}
\usetikzlibrary{arrows.meta,babel}
\usepackage[colorlinks=true,urlcolor=blue!50!black,linkcolor=blue!50!black]{hyperref}

\hyphenpenalty=10000 \exhyphenpenalty=10000 \emergencystretch=3em

% Colores de las figuras (los mismos que usa la página)
\definecolor{extremo}{RGB}{47,107,79}
\definecolor{azul}{RGB}{44,74,154}

% Definiciones numeradas por unidad
\theoremstyle{definition}
\newtheorem{defn}{Definición}
% Numeración por unidad: \unidad{u}{n} (después de \input{preambulo}) hace que las
% definiciones de la nota se llamen u.n, u.(n+1), ...; n es el primer número
% libre de esa unidad (ver la tabla de CONVENCIONES.md).
\newcommand{\unidad}[2]{\renewcommand{\thedefn}{#1.\arabic{defn}}\setcounter{defn}{\numexpr#2-1\relax}}

% Observaciones y erratas del video. Usar SOLO si hay algo que señalar:
%   \begin{erratas}
%     \errata{203}{3:23}{Se dice ..., pero debe decir ...}
%   \end{erratas}
% El primer argumento son los segundos (enlaza al minuto exacto del video),
% el segundo es la marca visible, el tercero es el texto.
\newenvironment{erratas}{\section*{Observaciones y erratas del video}\begin{itemize}}{\end{itemize}}
\newcommand{\errata}[3]{%
  \edef\errurl{https://www.youtube.com/watch?v=\videoid\string&t=#1}%
  \item[\textbf{\expandafter\href\expandafter{\errurl}{#2}}] #3}
) hace que las
% definiciones de la nota se llamen u.n, u.(n+1), ...; n es el primer número
% libre de esa unidad (ver la tabla de CONVENCIONES.md).
\newcommand{\unidad}[2]{\renewcommand{\thedefn}{#1.\arabic{defn}}\setcounter{defn}{\numexpr#2-1\relax}}

% Observaciones y erratas del video. Usar SOLO si hay algo que señalar:
%   \begin{erratas}
%     \errata{203}{3:23}{Se dice ..., pero debe decir ...}
%   \end{erratas}
% El primer argumento son los segundos (enlaza al minuto exacto del video),
% el segundo es la marca visible, el tercero es el texto.
\newenvironment{erratas}{\section*{Observaciones y erratas del video}\begin{itemize}}{\end{itemize}}
\newcommand{\errata}[3]{%
  \edef\errurl{https://www.youtube.com/watch?v=\videoid\string&t=#1}%
  \item[\textbf{\expandafter\href\expandafter{\errurl}{#2}}] #3}
.
% Compilar con LuaLaTeX:  lualatex nombre-de-la-nota.tex
\usepackage[spanish,provide=*]{babel}
\usepackage[T1]{fontenc}
\usepackage{lmodern}
\usepackage[a4paper,margin=2.6cm]{geometry}
\usepackage{amsmath,amssymb,amsthm}
\usepackage{xcolor}
\usepackage{tikz}
\usetikzlibrary{arrows.meta,babel}
\usepackage[colorlinks=true,urlcolor=blue!50!black,linkcolor=blue!50!black]{hyperref}

\hyphenpenalty=10000 \exhyphenpenalty=10000 \emergencystretch=3em

% Colores de las figuras (los mismos que usa la página)
\definecolor{extremo}{RGB}{47,107,79}
\definecolor{azul}{RGB}{44,74,154}

% Definiciones numeradas por unidad
\theoremstyle{definition}
\newtheorem{defn}{Definición}
% Numeración por unidad: \unidad{u}{n} (después de % Preámbulo compartido de las notas del curso de topología algebraica.
% Cada nota empieza con \documentclass[11pt]{article}, define \videoid{...}
% (el identificador del video de YouTube) y luego hace % Preámbulo compartido de las notas del curso de topología algebraica.
% Cada nota empieza con \documentclass[11pt]{article}, define \videoid{...}
% (el identificador del video de YouTube) y luego hace \input{preambulo}.
% Compilar con LuaLaTeX:  lualatex nombre-de-la-nota.tex
\usepackage[spanish,provide=*]{babel}
\usepackage[T1]{fontenc}
\usepackage{lmodern}
\usepackage[a4paper,margin=2.6cm]{geometry}
\usepackage{amsmath,amssymb,amsthm}
\usepackage{xcolor}
\usepackage{tikz}
\usetikzlibrary{arrows.meta,babel}
\usepackage[colorlinks=true,urlcolor=blue!50!black,linkcolor=blue!50!black]{hyperref}

\hyphenpenalty=10000 \exhyphenpenalty=10000 \emergencystretch=3em

% Colores de las figuras (los mismos que usa la página)
\definecolor{extremo}{RGB}{47,107,79}
\definecolor{azul}{RGB}{44,74,154}

% Definiciones numeradas por unidad
\theoremstyle{definition}
\newtheorem{defn}{Definición}
% Numeración por unidad: \unidad{u}{n} (después de \input{preambulo}) hace que las
% definiciones de la nota se llamen u.n, u.(n+1), ...; n es el primer número
% libre de esa unidad (ver la tabla de CONVENCIONES.md).
\newcommand{\unidad}[2]{\renewcommand{\thedefn}{#1.\arabic{defn}}\setcounter{defn}{\numexpr#2-1\relax}}

% Observaciones y erratas del video. Usar SOLO si hay algo que señalar:
%   \begin{erratas}
%     \errata{203}{3:23}{Se dice ..., pero debe decir ...}
%   \end{erratas}
% El primer argumento son los segundos (enlaza al minuto exacto del video),
% el segundo es la marca visible, el tercero es el texto.
\newenvironment{erratas}{\section*{Observaciones y erratas del video}\begin{itemize}}{\end{itemize}}
\newcommand{\errata}[3]{%
  \edef\errurl{https://www.youtube.com/watch?v=\videoid\string&t=#1}%
  \item[\textbf{\expandafter\href\expandafter{\errurl}{#2}}] #3}
.
% Compilar con LuaLaTeX:  lualatex nombre-de-la-nota.tex
\usepackage[spanish,provide=*]{babel}
\usepackage[T1]{fontenc}
\usepackage{lmodern}
\usepackage[a4paper,margin=2.6cm]{geometry}
\usepackage{amsmath,amssymb,amsthm}
\usepackage{xcolor}
\usepackage{tikz}
\usetikzlibrary{arrows.meta,babel}
\usepackage[colorlinks=true,urlcolor=blue!50!black,linkcolor=blue!50!black]{hyperref}

\hyphenpenalty=10000 \exhyphenpenalty=10000 \emergencystretch=3em

% Colores de las figuras (los mismos que usa la página)
\definecolor{extremo}{RGB}{47,107,79}
\definecolor{azul}{RGB}{44,74,154}

% Definiciones numeradas por unidad
\theoremstyle{definition}
\newtheorem{defn}{Definición}
% Numeración por unidad: \unidad{u}{n} (después de % Preámbulo compartido de las notas del curso de topología algebraica.
% Cada nota empieza con \documentclass[11pt]{article}, define \videoid{...}
% (el identificador del video de YouTube) y luego hace \input{preambulo}.
% Compilar con LuaLaTeX:  lualatex nombre-de-la-nota.tex
\usepackage[spanish,provide=*]{babel}
\usepackage[T1]{fontenc}
\usepackage{lmodern}
\usepackage[a4paper,margin=2.6cm]{geometry}
\usepackage{amsmath,amssymb,amsthm}
\usepackage{xcolor}
\usepackage{tikz}
\usetikzlibrary{arrows.meta,babel}
\usepackage[colorlinks=true,urlcolor=blue!50!black,linkcolor=blue!50!black]{hyperref}

\hyphenpenalty=10000 \exhyphenpenalty=10000 \emergencystretch=3em

% Colores de las figuras (los mismos que usa la página)
\definecolor{extremo}{RGB}{47,107,79}
\definecolor{azul}{RGB}{44,74,154}

% Definiciones numeradas por unidad
\theoremstyle{definition}
\newtheorem{defn}{Definición}
% Numeración por unidad: \unidad{u}{n} (después de \input{preambulo}) hace que las
% definiciones de la nota se llamen u.n, u.(n+1), ...; n es el primer número
% libre de esa unidad (ver la tabla de CONVENCIONES.md).
\newcommand{\unidad}[2]{\renewcommand{\thedefn}{#1.\arabic{defn}}\setcounter{defn}{\numexpr#2-1\relax}}

% Observaciones y erratas del video. Usar SOLO si hay algo que señalar:
%   \begin{erratas}
%     \errata{203}{3:23}{Se dice ..., pero debe decir ...}
%   \end{erratas}
% El primer argumento son los segundos (enlaza al minuto exacto del video),
% el segundo es la marca visible, el tercero es el texto.
\newenvironment{erratas}{\section*{Observaciones y erratas del video}\begin{itemize}}{\end{itemize}}
\newcommand{\errata}[3]{%
  \edef\errurl{https://www.youtube.com/watch?v=\videoid\string&t=#1}%
  \item[\textbf{\expandafter\href\expandafter{\errurl}{#2}}] #3}
) hace que las
% definiciones de la nota se llamen u.n, u.(n+1), ...; n es el primer número
% libre de esa unidad (ver la tabla de CONVENCIONES.md).
\newcommand{\unidad}[2]{\renewcommand{\thedefn}{#1.\arabic{defn}}\setcounter{defn}{\numexpr#2-1\relax}}

% Observaciones y erratas del video. Usar SOLO si hay algo que señalar:
%   \begin{erratas}
%     \errata{203}{3:23}{Se dice ..., pero debe decir ...}
%   \end{erratas}
% El primer argumento son los segundos (enlaza al minuto exacto del video),
% el segundo es la marca visible, el tercero es el texto.
\newenvironment{erratas}{\section*{Observaciones y erratas del video}\begin{itemize}}{\end{itemize}}
\newcommand{\errata}[3]{%
  \edef\errurl{https://www.youtube.com/watch?v=\videoid\string&t=#1}%
  \item[\textbf{\expandafter\href\expandafter{\errurl}{#2}}] #3}
) hace que las
% definiciones de la nota se llamen u.n, u.(n+1), ...; n es el primer número
% libre de esa unidad (ver la tabla de CONVENCIONES.md).
\newcommand{\unidad}[2]{\renewcommand{\thedefn}{#1.\arabic{defn}}\setcounter{defn}{\numexpr#2-1\relax}}

% Observaciones y erratas del video. Usar SOLO si hay algo que señalar:
%   \begin{erratas}
%     \errata{203}{3:23}{Se dice ..., pero debe decir ...}
%   \end{erratas}
% El primer argumento son los segundos (enlaza al minuto exacto del video),
% el segundo es la marca visible, el tercero es el texto.
\newenvironment{erratas}{\section*{Observaciones y erratas del video}\begin{itemize}}{\end{itemize}}
\newcommand{\errata}[3]{%
  \edef\errurl{https://www.youtube.com/watch?v=\videoid\string&t=#1}%
  \item[\textbf{\expandafter\href\expandafter{\errurl}{#2}}] #3}

\unidad{3}{23} % unidad 3; primer enunciado: 3.23

\title{El teorema de clasificación de cubrientes (unicidad)}
\author{Luis Jorge Sánchez Saldaña\\[2pt]\small Facultad de Ciencias, UNAM}
\date{}

\begin{document}
\maketitle

\noindent\small Notas del video
\href{https://www.youtube.com/watch?v=EZlgC_DT7D0}{\guillemotleft{}El teorema de clasificación de cubrientes - parte 2\guillemotright{}}
(curso de Introducción a la Topología Algebraica). Se define isomorfismo de cubrientes y se demuestra que dos cubrientes son isomorfos (respetando puntos base) si y sólo si tienen la misma imagen en el grupo fundamental.
\normalsize

\section*{Isomorfismos de cubrientes}

En el video anterior se demostró que, dado un subgrupo del grupo fundamental de $X$ (con todas las hipótesis usuales), siempre existe un cubriente que realiza ese subgrupo (Proposición~3.22). Ese es un problema de existencia, y la pregunta natural ahora es si hay algún tipo de unicidad. El objetivo de este video es responder afirmativamente esta pregunta; pero para hablar de unicidad primero se necesita una noción de isomorfismo entre cubrientes.

\begin{defn}[Isomorfismo de cubrientes]\label{def:iso-cubrientes}
Dos cubrientes $p_1\colon\tilde X_1\to X$ y $p_2\colon\tilde X_2\to X$ son \emph{isomorfos} si existe un homeomorfismo $f\colon\tilde X_1\to\tilde X_2$ tal que el diagrama
\[
  \begin{array}{ccc}
    \tilde X_1 & \xrightarrow{\ f\ } & \tilde X_2\\[2pt]
    & {\scriptstyle p_1}\searrow\quad\swarrow{\scriptstyle p_2} & \\[2pt]
    & X &
  \end{array}
\]
conmuta, es decir, $p_2\circ f=p_1$.
\end{defn}

Claramente la relación de isomorfismo entre cubrientes de un espacio base fijo es una relación de equivalencia. Que el diagrama conmute dice varias cosas. Por ejemplo, para cualquier punto $x\in X$, la función $f$ se restringe a una biyección
\[
  p_1^{-1}(x)\to p_2^{-1}(x).
\]
De hecho, si se tiene una vecindad regular $U$ en $X$, se puede levantar a $\tilde X_1$ y a $\tilde X_2$, y pasa lo obvio: las hojas sobre $U$ de un lado y del otro van a estar en biyección, y van a ser homeomorfas vía $f$. Así que un isomorfismo entre cubrientes dice que los cubrientes tienen esencialmente la misma estructura como cubrientes.

\section*{La unicidad}

\begin{prop}\label{prop:unicidad-cubrientes}
Si $X$ es arco-conexo y localmente arco-conexo, entonces $p_1\colon\tilde X_1\to X$ y $p_2\colon\tilde X_2\to X$ son isomorfos vía un homeomorfismo $f\colon\tilde X_1\to\tilde X_2$ con $f(\tilde x_1)=\tilde x_2$, donde $p_1(\tilde x_1)=p_2(\tilde x_2)=x_0$, si y sólo si
\[
  p_{1*}\pi_1(\tilde X_1,\tilde x_1)=p_{2*}\pi_1(\tilde X_2,\tilde x_2).
\]
\end{prop}

Recuérdese que ambas imágenes viven en $\pi_1(X,x_0)$. Esto da una forma de caracterizar el hecho de que dos cubrientes sean isomorfos, y es justamente un enunciado del tipo unicidad: dado un subgrupo $H$ del grupo fundamental, existe a lo más un cubriente, salvo isomorfismo, que lo realiza en el sentido de la Proposición~3.22.

\begin{proof}
\emph{($\Rightarrow$)} Del hecho de que se tenga el diagrama
\[
  \begin{array}{ccc}
    \tilde X_1 & \xrightarrow{\ f\ } & \tilde X_2\\[2pt]
    & {\scriptstyle p_1}\searrow\quad\swarrow{\scriptstyle p_2} & \\[2pt]
    & X &
  \end{array}
\]
y de que conmute, y de que $f$ induce un isomorfismo en grupos fundamentales, es bastante claro que las imágenes de los grupos fundamentales bajo $p_1$ y $p_2$ son iguales; no hay mucho que demostrar.

\medskip\noindent\emph{($\Leftarrow$)} Ahora se supone la igualdad de las imágenes y se piensa en términos del mismo diagrama. Ni siquiera hace falta que sean iguales: que la imagen bajo $p_{1*}$ de $\pi_1(\tilde X_1)$ esté contenida en la imagen bajo $p_{2*}$ de $\pi_1(\tilde X_2)$ dice, por el criterio de levantamiento (Proposición~3.11), que existe un levantamiento
\[
  \tilde p_1\colon(\tilde X_1,\tilde x_1)\to(\tilde X_2,\tilde x_2)
\]
de $p_1$ al cubriente $p_2$. De manera análoga, pensando ahora en el diagrama transpuesto y usando la otra inclusión, se obtiene un levantamiento
\[
  \tilde p_2\colon(\tilde X_2,\tilde x_2)\to(\tilde X_1,\tilde x_1)
\]
de $p_2$ al cubriente $p_1$. Así se obtiene el diagrama
\[
  \begin{array}{ccccc}
    & & \tilde X_1 & & \\[2pt]
    & {\scriptstyle\tilde p_2}\nearrow & \big\downarrow{\scriptstyle p_1} & \searrow{\scriptstyle\tilde p_1} & \\[2pt]
    \tilde X_2 & \xrightarrow{\ p_2\ } & X & \xleftarrow{\ p_2\ } & \tilde X_2
  \end{array}
\]
La composición $\tilde p_1\circ\tilde p_2\colon\tilde X_2\to\tilde X_2$ es un levantamiento de $p_2$ al cubriente $p_2$ que manda $\tilde x_2$ en $\tilde x_2$; pero la identidad también lo es. Como también se demostró que dicho levantamiento es único (Proposición~3.12), a este levantamiento no le queda otra que ser la identidad:
\[
  \tilde p_1\circ\tilde p_2=\mathrm{Id}_{\tilde X_2}.
\]
Análogamente, $\tilde p_2\circ\tilde p_1=\mathrm{Id}_{\tilde X_1}$. Por lo tanto $\tilde p_1$ es un homeomorfismo con $p_2\circ\tilde p_1=p_1$ y $\tilde p_1(\tilde x_1)=\tilde x_2$, y los dos cubrientes son isomorfos.
\end{proof}

Esta demostración fue bastante más sencilla que las anteriores, pero aun así el resultado no deja de ser poderoso: la realización en cubrientes de cualquier subgrupo del grupo fundamental de $X$ es única.

\begin{erratas}
  \errata{328}{5:28}{En el pizarrón se escribe dos veces $\tilde p_2\circ\tilde p_1$: primero $\tilde p_2\circ\tilde p_1=\mathrm{Id}_{\tilde X_2}$ y luego $\tilde p_2\circ\tilde p_1=\mathrm{Id}_{\tilde X_1}$ (y de palabra se dice lo mismo). Como $\tilde p_2\circ\tilde p_1$ va de $\tilde X_1$ a $\tilde X_1$, la primera debe decir $\tilde p_1\circ\tilde p_2=\mathrm{Id}_{\tilde X_2}$.}
\end{erratas}

\end{document}
